% Source: https://artofproblemsolving.com/community/c3h1623534_proving_a_triangle_theorem
% Stable identity: aops-topic-1623534-q01
Measure lengths in units of the small side. Label the grid points by $(i,j)$, where $i=0,1,\dots,n$ is the row of points counted from the top vertex and $0\le j\le i$ is the position in the row; row $i$ has $i+1$ points. Every triangle in the figure has sides along the grid lines and is either pointing up or pointing down.

\textbf{Upward triangles.} An upward triangle of side $s$ is fixed by its top vertex $(i,j)$, with other vertices $(i+s,j)$ and $(i+s,j+s)$; it fits exactly when $i+s\le n$. The number of possible top vertices is the number of points in rows $0$ to $n-s$, namely $1+2+\cdots+(n-s+1)=\dbinom{n-s+2}{2}$. Summing over $s=1,\dots,n$ and using the hockey-stick identity,
\[
U=\sum_{s=1}^{n}\dbinom{n-s+2}{2}=\dbinom22+\dbinom32+\cdots+\dbinom{n+1}{2}=\dbinom{n+2}{3}=\dfrac{n(n+1)(n+2)}{6}.
\]

\textbf{Downward triangles.} A downward triangle of side $s$ has top vertices $(i,j)$ and $(i,j+s)$ and bottom vertex $(i+s,j+s)$; it fits when $j+s\le i$ and $i+s\le n$. For each row $i$ from $s$ to $n-s$ there are $i-s+1$ choices of $j$, so the number is $1+2+\cdots+(n-2s+1)=\dbinom{n-2s+2}{2}$, for $1\le s\le\dfrac {n}{2}$.

\textbf{(a)} For $n=5$: $U=\dbinom73=35$ and the downward triangles number $\dbinom52+\dbinom32=10+3=13$, so there are $35+13=\boxed{48}$ triangles. For $n=7$: $U=\dbinom93=84$ and the downward triangles number $\dbinom72+\dbinom52+\dbinom32=21+10+3=34$, so there are $84+34=\boxed{118}$ triangles.

\textbf{(b)} Let $n=2k+1$ with $k\ge1$. The downward triangles have sides $s=1,\dots,k$, and with $m=k-s+1$ (so $n-2s+2=2m+1$, and $m$ runs from $1$ to $k$) their number is
\[
D=\sum_{m=1}^{k}\dbinom{2m+1}{2}=\sum_{m=1}^{k}m(2m+1)=2\cdot\dfrac{k(k+1)(2k+1)}{6}+\dfrac{k(k+1)}{2}=\dfrac{k(k+1)(4k+5)}{6}.
\]
Also $U=\dfrac{(2k+1)(2k+2)(2k+3)}{6}=\dfrac{(k+1)(2k+1)(2k+3)}{3}$. Hence the total is
\[
U+D=\dfrac{(k+1)\left[2(2k+1)(2k+3)+k(4k+5)\right]}{6}=\dfrac{(k+1)(12k^2+21k+6)}{6}=\dfrac{(k+1)(4k^2+7k+2)}{2}=\dfrac{4k^3+11k^2+9k+2}{2}.
\]
On the other hand $(2k+1)(2k+3)(4k+3)=16k^3+44k^2+36k+9$, so
\[
\dfrac{n(n+2)(2n+1)-1}{8}=\dfrac{16k^3+44k^2+36k+8}{8}=\dfrac{4k^3+11k^2+9k+2}{2}.
\]
The two expressions agree, which proves the formula for odd $n$ (for $n=1$ it gives $1$, also correct). $\blacksquare$

(For even $n=2k$ the same method gives $D=\sum_{m=1}^{k}\dbinom{2m}{2}=\dfrac{k(k+1)(4k-1)}{6}$ and the total $\dfrac{n(n+2)(2n+1)}{8}$, the formula quoted in the post. Both cases together can be written as $\left\lfloor\dfrac{n(n+2)(2n+1)}{8}\right\rfloor$.)
